\documentclass[11pt]{article}
\usepackage[margin=1in]{geometry}
\usepackage{amsmath,amssymb,amsthm,mathtools}
\usepackage{graphicx}
\usepackage[hidelinks]{hyperref}
\usepackage{xcolor}
\usepackage{microtype}

\newtheorem{theorem}{Theorem}
\newtheorem{lemma}[theorem]{Lemma}
\newtheorem{corollary}[theorem]{Corollary}
\newtheorem{remark}[theorem]{Remark}
\newcommand{\R}{\mathbb R}
\newcommand{\E}{\mathcal E}

\title{The Converse for Spherically Symmetric Trees\\
\large A candidate full solution to the question after Theorem 2.2 of arXiv:1008.3418}
\author{Automated mathematics run \texttt{fa\_banach\_001}, lane 9}
\date{27 August 2026}

\begin{document}
\maketitle

\begin{center}
\fbox{\parbox{0.91\textwidth}{\textbf{Status.}
Candidate full solution, likely valid, pending expert review.  The argument is
self-contained apart from the standard equivalence between generalized
roundness and supremal negative type recalled in the source paper.}}
\end{center}

\section{Source question and answer}

Caffarelli, Doust, and Weston prove in Theorem 2.2 of
\cite{CaffarelliDoustWeston} that if the downward degree sequence
$(d_0,d_1,\ldots)$ of a countable spherically symmetric tree (SST) has
infinitely many entries larger than one, then the tree has generalized
roundness one.  On page 5, immediately before Section 3, they ask:

\begin{quote}
``We do not know if the converse of Theorem 2.2 is true.''
\end{quote}

\begin{figure}[ht]
  \centering
  \includegraphics[width=\textwidth]{figures/open_problem_crop.png}
  \caption{Theorem 2.2 and the converse question on page 5 of the source paper.}
\end{figure}

The answer is affirmative.  We prove a stronger metric-tree statement.  A
\emph{finite-topological tree} will mean a geodesic metric tree obtained by
gluing finitely many compact intervals and half-lines at endpoints, without
creating a cycle.  Its open edges are the interval or half-line interiors
after a finite vertex set has been chosen.  Subdivision only increases the
number of open edges, so no minimality convention is needed.

\section{Proof intuition}

For a finite signed configuration of total mass zero on a tree, there is a
unique cumulative flow along each oriented edge.  Integrating the
distance-power form twice transfers the coefficients to these flows.  On one
edge the mixed derivative of $|s-t|^p$ has the favorable sign and produces the
positive Riesz kernel $|s-t|^{p-2}$.  Between two distinct edges, the distance
is affine in the two edge coordinates, so the unfavorable contribution has
the shifted Hankel kernel $(s+t+c)^{p-2}$.

The key point is quantitative: when $p$ is just above one, every Hankel cross
term is a small fraction of the corresponding within-edge Riesz energies.  A
logarithmic change of variables converts both homogeneous kernels into
convolution kernels.  Their Fourier-multiplier ratio is exactly
\[
  \frac{\cos(\pi\alpha/2)}{\cosh(\pi\xi)},\qquad \alpha=2-p.
\]
Hence the relative norm is $\cos(\pi\alpha/2)
=\sin(\pi(p-1)/2)$, which tends to zero as $p\downarrow1$.
There are only finitely many edge pairs, so diagonal dominance proves
$p$-negative type for some $p>1$.

\section{The Riesz--Hankel estimate}

Fix $0<\alpha<1$.  For bounded compactly supported real functions on
$\R_+=(0,\infty)$, set
\begin{align*}
 A_\alpha(f,g)
   &=\int_0^\infty\!\int_0^\infty
       f(s)g(t)|s-t|^{-\alpha}\,ds\,dt,\\
 H_{\alpha,c}(f,g)
   &=\int_0^\infty\!\int_0^\infty
       f(s)g(t)(s+t+c)^{-\alpha}\,ds\,dt
       \qquad(c\ge0).
\end{align*}
Write $A_\alpha(f)=A_\alpha(f,f)$.

\begin{lemma}[Sharp half-line comparison]\label{lem:hankel}
For all such $f,g$ and every $c\ge0$,
\[
 |H_{\alpha,c}(f,g)|
 \le r_\alpha A_\alpha(f)^{1/2}A_\alpha(g)^{1/2},
 \qquad r_\alpha=\cos\frac{\pi\alpha}{2}.
\]
In particular, $A_\alpha$ is a positive definite quadratic form.
\end{lemma}

\begin{proof}
We first take $c=0$.  Make the logarithmic substitution $s=e^x$ and define
\[
  (Uf)(x)=e^{(1-\alpha/2)x}f(e^x).
\]
Direct cancellation of the powers of $e^x$ and $e^y$ gives
\begin{align*}
 A_\alpha(f,g)
 &=\int_{\R^2}(Uf)(x)(Ug)(y)a_\alpha(x-y)\,dx\,dy,\\
 H_{\alpha,0}(f,g)
 &=\int_{\R^2}(Uf)(x)(Ug)(y)b_\alpha(x-y)\,dx\,dy,
\end{align*}
where
\[
 a_\alpha(u)=\bigl(2|\sinh(u/2)|\bigr)^{-\alpha},
 \qquad
 b_\alpha(u)=\bigl(2\cosh(u/2)\bigr)^{-\alpha}.
\]
Both kernels are integrable because $0<\alpha<1$.  With any consistent
Fourier-transform convention, beta integration gives, for
$z=\alpha/2+i\xi$,
\begin{align*}
 \widehat b_\alpha(\xi)
 &=\frac{|\Gamma(z)|^2}{\Gamma(\alpha)},\\
 \widehat a_\alpha(\xi)
 &=\Gamma(1-\alpha)
 \left(
 \frac{\Gamma(z)}{\Gamma(1-\overline z)}
 +\frac{\Gamma(\overline z)}{\Gamma(1-z)}
 \right).
\end{align*}
The reflection identity for the gamma function yields
\[
 \widehat a_\alpha(\xi)
 =\frac{2\Gamma(1-\alpha)|\Gamma(z)|^2
          \sin(\pi\alpha/2)\cosh(\pi\xi)}{\pi}>0.
\]
Using $\Gamma(\alpha)\Gamma(1-\alpha)=\pi/\sin(\pi\alpha)$,
we therefore obtain the exact quotient
\begin{equation}\label{eq:ratio}
 \frac{\widehat b_\alpha(\xi)}{\widehat a_\alpha(\xi)}
 =\frac{\cos(\pi\alpha/2)}{\cosh(\pi\xi)}
 \le r_\alpha.
\end{equation}
Plancherel and weighted Cauchy--Schwarz now prove the asserted bilinear
estimate for $c=0$, while positivity of $\widehat a_\alpha$ proves positivity
of $A_\alpha$.

For $c>0$, translate both zero-extended functions to the right by $c/2$:
\[
 f_c(u)=f(u-c/2)\mathbf 1_{(c/2,\infty)}(u),
 \qquad
 g_c(u)=g(u-c/2)\mathbf 1_{(c/2,\infty)}(u).
\]
Then $H_{\alpha,c}(f,g)=H_{\alpha,0}(f_c,g_c)$, whereas translation
invariance of $|s-t|^{-\alpha}$ gives
$A_\alpha(f_c)=A_\alpha(f)$ and similarly for $g$.  The result follows.
\end{proof}

The same estimate applies when either function is supported on a finite
interval: extend it by zero to the half-line.  Reversing an interval coordinate
does not change its diagonal Riesz energy.

\section{Flows on a finite-topological tree}

Let $G$ be a finite-topological tree with open edges
$e=1,\ldots,N$.  Choose a root and orient each edge away from it.  For a
finitely supported signed measure
\[
  \mu=\sum_{k=1}^m a_k\delta_{x_k},\qquad \sum_{k=1}^m a_k=0,
\]
define $F_e(s)$ to be the $\mu$-mass of the component lying beyond the point
at coordinate $s$ on edge $e$.  On every half-line $F_e$ has compact support.
At a branch vertex, the incoming value minus the sum of the outgoing values
is exactly the atom of $\mu$ at that vertex.  Consequently, edgewise
integration by parts gives
\begin{equation}\label{eq:flowibp}
  \int_G \varphi\,d\mu
  =\sum_{e=1}^N\int_e F_e(s)\,\partial_e\varphi(s)\,ds
\end{equation}
for every continuous function that is continuously differentiable on open
edges.  This identity incorporates atoms at vertices through the preceding
Kirchhoff balance; there are no omitted boundary terms.

\begin{lemma}[Finite-edge coercivity]\label{lem:tree}
Let $1<p<2$, put $\alpha=2-p$, and let
$r_p=\sin(\pi(p-1)/2)=\cos(\pi\alpha/2)$.  Then
\begin{equation}\label{eq:coercive}
 -\int_G\!\int_G d(x,y)^p\,d\mu(x)d\mu(y)
 \ge p(p-1)\bigl(1-(N-1)r_p\bigr)
       \sum_{e=1}^N A_\alpha(F_e).
\end{equation}
In particular, if $(N-1)r_p\le1$, then $G$ has $p$-negative type.
\end{lemma}

\begin{proof}
Apply \eqref{eq:flowibp} in both variables.  Since
$\alpha\in(0,1)$, all resulting second-derivative kernels are locally
integrable, so the calculation is justified by ordinary approximation or
Fubini's theorem.  On a single edge,
\[
 -\partial_s\partial_t|s-t|^p
 =p(p-1)|s-t|^{-\alpha}.
\]
For two distinct open edges $e\ne f$, the path joining their interior points
exits each edge through a fixed endpoint.  If $u,v$ denote distances from
those endpoints, then
\[
 d(x_e(u),x_f(v))=u+v+c_{ef}
\]
for some $c_{ef}\ge0$.  Relative to the globally chosen orientations,
$u$ or $v$ may be reversed; hence the corresponding mixed derivative has an
inessential sign.  It follows that
\begin{align*}
 -\int_G\!\int_G d(x,y)^p\,d\mu(x)d\mu(y)
 =p(p-1)\bigg(&\sum_e A_\alpha(F_e)\\
 &+\sum_{e\ne f}\varepsilon_{ef}
 H_{\alpha,c_{ef}}(\widetilde F_e,\widetilde F_f)\bigg),
\end{align*}
where $\varepsilon_{ef}\in\{-1,1\}$ and each tilde records at most a
coordinate reversal and zero extension.

Lemma \ref{lem:hankel} bounds the absolute value of the $(e,f)$ term by
$r_p A_\alpha(F_e)^{1/2}A_\alpha(F_f)^{1/2}$.  Therefore
\begin{align*}
 \frac{-\iint d(x,y)^p\,d\mu(x)d\mu(y)}{p(p-1)}
 &\ge \sum_e A_\alpha(F_e)
 -r_p\sum_{e\ne f}\sqrt{A_\alpha(F_e)A_\alpha(F_f)}\\
 &\ge \bigl(1-(N-1)r_p\bigr)\sum_e A_\alpha(F_e).
\end{align*}
The last inequality follows from
$(\sum_e\sqrt{A_e})^2\le N\sum_e A_e$.  This proves
\eqref{eq:coercive}.  When its coefficient is nonnegative, the defining
zero-sum quadratic inequality for $p$-negative type follows.
\end{proof}

\begin{theorem}[Finite-topological trees have roundness above one]
\label{thm:finite-topological}
Every finite-topological tree has generalized roundness strictly greater than
one.  More explicitly, if it is presented with $N$ open edges, then it has
$p$-negative type for every $1<p<2$ satisfying
\begin{equation}\label{eq:range}
 (N-1)\sin\frac{\pi(p-1)}2\le1.
\end{equation}
For example, the universally safe choice
\[
 p=1+\frac2\pi\arcsin\!\left(\frac{1}{2\max\{1,N-1\}}\right)
\]
is strictly larger than one.
\end{theorem}

\begin{proof}
Lemma \ref{lem:tree} proves the zero-sum quadratic inequality for every finite
configuration in $G$.  This is $p$-negative type.  Generalized roundness equals
supremal negative type, as recalled in the source paper
\cite{CaffarelliDoustWeston}, so the displayed safe exponent gives
$\gamma(G)>1$.
\end{proof}

\section{Resolution of the converse}

\begin{corollary}[Converse of source Theorem 2.2]
Let $(T,\rho)$ be a countable spherically symmetric tree with downward degree
sequence $(d_0,d_1,\ldots)$.  Then
\[
 \gamma(T,\rho)=1
 \quad\Longleftrightarrow\quad
 \bigl|\{j:d_j>1\}\bigr|=\aleph_0.
\]
\end{corollary}

\begin{proof}
The reverse implication is Theorem 2.2 of
\cite{CaffarelliDoustWeston}.  Conversely, suppose that only finitely many
$d_j$ exceed one.  Because the tree is countably infinite, its downward
degrees are positive; after the last branching level every descendant chain
is a ray.  There are finitely many such rays and only a finite subtree before
them.  Suppressing degree-two vertices in its geometric realization therefore
presents it as a finite-topological tree with some finite number $N$ of open
edges.  The original vertex set $T$ is an isometric subset of that realization.
By Theorem \ref{thm:finite-topological}, the realization, and hence its subset
$T$, has $p$-negative type for some $p>1$.  Thus $\gamma(T)>1$.
\end{proof}

\section{Verification and limitations}

The proof gives a lower bound depending on the chosen finite edge
decomposition, not the exact generalized roundness.  Subdividing edges makes
the bound weaker but cannot invalidate it.  The argument also proves the
stronger result for arbitrary finite-topological trees, without spherical
symmetry or unit edge lengths.

The accompanying script \texttt{code/verify\_finite\_edge\_bound.py} checks
the closed multiplier ratio on five $\alpha$ values and 4,001 frequencies per
value, and independently compares the beta/gamma formulas with that closed
ratio at 25 parameter pairs.  It also tests the explicit safe exponent on
2,400 random weighted finite trees by compressing the powered distance matrix
to the zero-sum subspace.  These checks are designed to catch sign,
factor-of-two, and edge-count errors; they are not used in the proof.

The points most deserving human verification are (i) the edgewise
integration-by-parts identity at branch vertices, and (ii) the fact that the
sum over distinct edges is an ordered sum, matching the factor
$N-1$ in the final Cauchy--Schwarz estimate.

\section{Bounded novelty check}

On 27 August 2026, the run's registry, solution, attempt, and proof-gap indexes
were searched for arXiv:1008.3418, its title, the exact converse wording, and
finite-ended/finitely-branching negative-type variants.  Bounded arXiv and web
searches used the same phrases and inspected the related arXiv papers
1608.03699, 0901.0695, and 1108.0451.  The searches found the original question
and adjacent results on generalized roundness and finite metric spaces, but no
exact answer to this converse and no matching finite-topological-tree theorem.
This makes novelty plausible but does not certify it.

\begin{thebibliography}{9}
\bibitem{CaffarelliDoustWeston}
E. Caffarelli, I. Doust, and A. Weston,
\emph{Metric trees of generalized roundness one},
Aequationes Math. 83 (2012), arXiv:1008.3418v2,
doi:10.1007/s00010-011-0108-8.
\end{thebibliography}

\end{document}
